\documentclass[11pt]{article}
\usepackage[a4paper,margin=25mm]{geometry}
\usepackage{amsmath,amssymb,mathtools,graphicx,booktabs,xcolor}
\usepackage[most]{tcolorbox}
\newcommand{\E}{\mathbb{E}}
\newcommand{\defeq}{\stackrel{\mathrm{def}}{=}}
\newcommand{\arginf}{\mathop{\mathrm{arg\,inf}}}
\newcommand{\norm}[1]{\lVert#1\rVert}
\newcommand{\IMFt}{\widetilde S}
\newcommand{\Wcal}{\mathcal W}
\newtcolorbox{idea}{colback=teal!6,colframe=teal!60!black,boxrule=0.6pt,arc=2pt,left=6pt,right=6pt,top=4pt,bottom=4pt}
\newtcolorbox{formula}{colback=orange!6,colframe=orange!70!black,boxrule=0.6pt,arc=2pt,left=6pt,right=6pt,top=4pt,bottom=4pt}
\setlength{\parskip}{4pt}\setlength{\parindent}{0pt}
\emergencystretch=3em\sloppy
\begin{document}
\begin{center}
{\LARGE Why later IRMF components stay accurate}\\[4pt]
{\large A guided tour of the contraction argument for components \(k\ge2\)}\\[4pt]
{\normalsize Companion to \texttt{irmf\_contraction\_note.pdf} (version 2). Work in progress.}
\end{center}

\begin{idea}
\textbf{What you will take away.}
\begin{enumerate}\setlength{\itemsep}{1pt}
\item Every IRMF stage subtracts a noisy fit, so errors could pile up.  They
  do not, and the reason is visible in frequency: each component takes its
  own band of frequencies.
\item In the sup norm (the largest error over all time points), one stage on
  its own can double an error.  A chain of stages cannot, because the
  stages work at well-separated scales.
\item The robust fit behaves, to first order, like a linear smoother with
  \emph{bent weights}, \(W_\phi\).  The Fisher expansion of the first
  component, applied at every stage, makes this exact.
\item What remains is a sum of small second-order terms.  It is largest at
  the finest windows under heavy contamination.
\end{enumerate}
\end{idea}

\paragraph{The test setting used in all figures.}
One smooth test signal on \(n=1000\) equally spaced points, periodic at the
ends: a slow baseline plus a chirp, with standard deviation 0.39.  Noise
\(\sigma=0.4\), contrast \(\rho_H\) with \(H=0.8\), weights \((1-|u|)_+^2\),
nine stages with windows of \(501,355,251,\dots,31\) points (\(a=\sqrt2\)).
Some figures also use the same signal multiplied by three, or add
contamination to 20\% of the points.


\section{The question}

IRMF fits a robust local level with a wide window, subtracts it, and repeats
on the residual with a narrower window (IMF.pdf, (3.2)):
\[
  \IMFt^{(k)}_t\defeq\arginf_x\sum_u\rho\bigl(Y^{(k)}_u-x\bigr)\,w^{(k)}_{u,t},
  \qquad Y^{(k+1)}\defeq Y^{(k)}-\IMFt^{(k)},\qquad h_{k+1}=h_k/a .
\]
The first component was handled by the Fisher expansion
(Proposition~4.1 of \texttt{proposition\_4\_1\_lemma\_4\_2.tex}):
\[
  \IMFt_t-S_t=F_t^{-1}\nabla_t+R_t,
\]
whose leading term has size \(\sqrt{z_x/(nh)}\).  The question is whether
component \(k\) keeps the error size of \emph{its own} window,
\(\sqrt{z_x/(nh_k)}\), or whether the errors of earlier stages pile up.

Why worry?  Stage \(k\) sees the error of stage \(k-1\) inside its input, and
subtracting a fit can amplify a signal.  For any smoother \(W\) with
nonnegative weights summing to one, the sup norm satisfies
\(\norm{I-W}_\infty=2-2\,W[t,t]\approx2\): a spike at one point, minus a
spread-out bump, has absolute mass 2.  Bounding one factor at a time gives
\(2^{k-1}\) after \(k\) stages.  The bound in
\texttt{second\_irmf\_error\_targets.tex} hits this wall: its constant
\(\Gamma_1\ge2\) is raised to the power \(k\).

\section{Warm-up: linear IMF in frequency language}

For \(\rho(x)=x^2/2\) everything is linear, and IMF.pdf, Lemma~2.3 gives
\[
  \IMFt^{(k)}-S^{(k)}=W^{(k)}A^{(k-1)}\varepsilon,\qquad
  A^{(k-1)}=(I-W^{(k-1)})\cdots(I-W^{(1)}).
\]
On the periodic grid every \(W^{(k)}\) is a convolution, so each acts on a
frequency \(\omega\) by multiplying it by a number \(\hat w_k(\omega)\), its
\emph{gain}.  Products of operators become products of gains:
\[
  \text{gain of component }k\text{'s error filter}
  =\hat w_k(\omega)\prod_{j<k}\bigl(1-\hat w_j(\omega)\bigr).
\]

\begin{figure}[h]\centering
\includegraphics[width=\textwidth]{T1_frequency.pdf}
\caption{\textbf{How linear IMF shares noise between stages.}
Each curve says how strongly a filter lets a frequency through: 1 means
fully, 0 means not at all.  Colour marks the stage, from dark (stage~1,
widest window) to light (stage~9, narrowest).
\textbf{Left:} one smoother lets slow variations through and blocks fast
ones; a narrower window lets more through.
\textbf{Middle:} the filter that produces the error of component \(k\).
Each one is a hump over its own range of frequencies.  The dashed curve is
what is left for the final residual.
\textbf{Right:} the subtraction \(I-W^{(k)}\) of each stage (one row per
stage).  Dark means the frequency is removed, light means it is kept.
\emph{Takeaway:} every frequency is handled by one or two stages, so noise is
shared out between stages, not multiplied.}
\end{figure}

Read panel b as a filter bank: every frequency is assigned mostly to one or
two components, and the bumps sum to at most one.  Noise therefore cannot
grow from stage to stage in this picture.  It is only \emph{shared out}
across stages.  For a positive-definite kernel every factor has gain in
\([0,1]\), so \(\norm{A^{(k)}}_2\le1\) in the \(\ell_2\) (energy) norm.

\section{The sup norm: one stage can double, a chain cannot}

Energy is not the sup norm.  Figure~\ref{fig:impulse} shows what happens to
the largest error.

\begin{figure}[h]\centering
\includegraphics[width=\textwidth]{T3_impulse_norms.pdf}
\caption{\textbf{How large the error filters are in the sup norm} (the largest
error over all time points).
\textbf{Left:} the filter that produces the error of component \(k\),
drawn in time.  The horizontal axis is measured in window half-widths of
that stage.  The curves for stages 3, 5, 7 and 9 lie on top of each other:
every stage uses the same filter shape at its own scale.
\textbf{Right:} three sizes per stage, on a log scale.  Orange: the worst
case if every stage could double the error, \(2^{k-1}\).  Purple: the actual
size of the chain of subtractions \(A^{(k-1)}\); it stays near 2.  Green: the
actual size of the error filter of component \(k\); it stays below 0.5 from
stage~2 on.
\emph{Takeaway:} the worst case never happens.}
\label{fig:impulse}
\end{figure}

\begin{idea}
\textbf{Why the chain does not grow: scale separation.}  Look again at
Figure~1c.  At any frequency, the coarse stages have already switched to
``keep'' (gain \(\approx1\)) and the fine stages are still at ``remove''
(gain \(\approx0\)).  Only one or two stages are in between.  A product of
many factors is therefore never a product of many numbers that are all
slightly above or below one.  It is a few ones, one transition, and zeros.
The proof (Lemma~2 of the note) turns this into a bound on the \(\ell_1\)
norm of the composite filter, using one fact: the finest factor
\(1-\hat w_{k-1}\) vanishes quadratically exactly where the coarser factors
are still changing.
\end{idea}

The resulting statement, uniform in \(k\), is
\[
  \norm{W^{(k)}A^{(k-1:l)}}_\infty\le C_*,\qquad \norm{A^{(m:l)}}_\infty\le1+C_*,
\]
for all partial products \(A^{(m:l)}=(I-W^{(m)})\cdots(I-W^{(l)})\).  On a
21-stage schedule (\(n=2^{16}\)) the computed values stay at most
\(1.00\) and \(2.14\).

\section{An aside: the Epanechnikov kernel}

IMF.pdf, Lemma~2.2 claims \(\norm{A^{(k)}}\le1\).  That fails for
Epanechnikov \(\tfrac34(1-u^2)_+\), because its Fourier transform has
negative lobes: \(1-\hat w\) can exceed one.  Figure~\ref{fig:epa} shows why
this does no real harm.

\begin{figure}[h]\centering
\includegraphics[width=\textwidth]{T4_epanechnikov.pdf}
\caption{\textbf{The Epanechnikov kernel.}
Each curve shows, for one stage, by how much the subtraction \(I-W^{(k)}\)
scales each frequency.  The vertical axis is zoomed near 1; a value above the
orange line means that frequency is amplified.
\textbf{Left}, kernel \((1-|u|)_+^2\): no stage amplifies anything.
\textbf{Right}, Epanechnikov: every stage amplifies a little, up to 8.7\%,
but each at different frequencies.
The dashed curve is the product over all nine stages.  It peaks at 1.077,
so the small amplifications do not add up.}
\label{fig:epa}
\end{figure}

The overshoot happens only where \(nh_k|\omega|\gtrsim4.5\), and because
the bandwidths shrink geometrically, at a fixed frequency the overshoots add
up to a bounded total.  So the kernel question is not fatal for the
contraction.  The bound needs Fourier decay of order \(\omega^{-2}\); it does
not need a third derivative of \(K\).

\section{From linear to robust}

\subsection{Making every stage a first-component problem}

For the robust fit, we need a reference to measure errors against.  Use the
\emph{\(\phi\)-path}: run the same recursion on the clean signal, with the
noise-averaged contrast \(\phi(x)=\E\rho(x+\varepsilon)\):
\[
  X^{(1)}_\phi\defeq X,\qquad
  S^{(k)}_{\phi,t}\defeq\arginf_x\sum_u\phi\bigl(X^{(k)}_{\phi,u}-x\bigr)w^{(k)}_{u,t},
  \qquad X^{(k+1)}_\phi\defeq X^{(k)}_\phi-S^{(k)}_\phi .
\]
Why this path: feed stage \(k\) the clean residual plus \emph{fresh} noise,
\(X^{(k)}_\phi+\varepsilon\).  That is exactly the first-component problem
with signal \(X^{(k)}_\phi\), and its population target is
\(S^{(k)}_\phi\).  So Proposition~4.1 applies at every stage, unchanged:
\[
  \text{single pass at stage }k:\qquad
  \mathcal M_k(X^{(k)}_\phi+\varepsilon)-S^{(k)}_\phi
  =F^{(k)\,-1}\nabla^{(k)}+R^{(k)} .
\]
The actual stage-\(k\) input differs from this by the errors made so far,
\(E^{(k)}=\sum_{j<k}e^{(j)}\): \(Y^{(k)}=X^{(k)}_\phi+\varepsilon-E^{(k)}\).

IMF.pdf, Section~3.2 defines the ``true IMFs'' \(S^{(k)}\) with the law of
the actual residual noise \(\varepsilon-E^{(k)}\) instead of
\(\varepsilon\).  Since \(E^{(k)}\) is small, the two targets differ only at
second order.  In 400 simulated repeats they were indistinguishable: the
difference was within Monte Carlo error at every stage.

\subsection{The robust smoother as a linear smoother with bent weights}

How does the stage-\(k\) fit respond to a small change \(E\) in its input?
To first order, through the derivative of the leading term
\(F^{-1}\nabla\), averaged over the noise:
\[
  W^{(k)}_\phi[t,u]\defeq\frac{\phi''\bigl(X^{(k)}_{\phi,u}-S^{(k)}_{\phi,t}\bigr)\,w^{(k)}_{u,t}}{F^{(k)}_t}.
\]
This is a deterministic matrix whose rows are nonnegative and sum to one, a
smoother like \(W^{(k)}\).  Its weights are the kernel weights multiplied by
the curvature \(\phi''\) at each point's distance from the local fit.

\begin{figure}[h]\centering
\includegraphics[width=\textwidth]{T5_kernel_tilt.pdf}
\caption{\textbf{What the robust fit listens to}, for one stage-1 window (centred at
the dotted line).
\textbf{Top:} the clean signal inside the window (black) and the value the
fit aims at (orange).
\textbf{Bottom:} how much each point counts.  Grey dashed: linear weights.
Green: weights of the linearised robust fit, \(W^{(1)}_\phi\).
\textbf{Left:} for the test signal the two weight curves almost coincide.
\textbf{Right:} with the signal three times larger, points whose value is
far from the orange line count less, and points close to it count more.
\emph{Takeaway:} to first order the robust fit is a linear smoother whose
weights bend towards points near the current fit.}
\label{fig:tilt}
\end{figure}

\begin{idea}
\textbf{Bent weights.}  The robust fit listens most to points close to its
current value.  If the signal is nearly flat on the scale of
\(B=\sqrt{H^2+\sigma^2}\), the bending cancels in the normalisation and
\(W_\phi=W\): robust IMF behaves like linear IMF to first order.  If the
signal varies more, the weights bend, but they stay a smoother.
\end{idea}

\section{The decomposition}

Put the two pieces together at every stage: fresh noise enters through the
Fisher leading term, and old errors are partly removed by \(W_\phi\), exactly
as in linear IMF.  Collecting the stages (Lemma~1 of the note) gives an
exact identity:

\begin{formula}
\[
  \underbrace{\IMFt^{(k)}-S^{(k)}_\phi}_{\text{error }e^{(k)}}
  =\underbrace{\sum_{j\le k}T^{(k,j)}F^{(j)\,-1}\nabla^{(j)}}_{\text{leading term }L^{(k)}}
  +\underbrace{\sum_{j\le k}T^{(k,j)}\bigl(R^{(j)}-\Phi^{(j)}\bigr)}_{\text{remainder}},
\]
\[
  T^{(k,k)}=I,\qquad T^{(k,j)}=-W^{(k)}_\phi(I-W^{(k-1)}_\phi)\cdots(I-W^{(j+1)}_\phi).
\]
\(R^{(j)}\): the Fisher remainder of the stage-\(j\) single pass
(Proposition~4.1).  \(\Phi^{(j)}\): the error of treating the response to
\(E^{(j)}\) as linear.
\end{formula}

Three remarks make the formula readable.
\begin{itemize}
\item In the linear case the remainder is zero and the leading term is IMF.pdf's
  \(W^{(k)}A^{(k-1)}\varepsilon\).
\item For IRMF the leading term is ``linear IMF with bent weights, acting on
  bounded scores''.  The transfer operators \(T^{(k,j)}\) are exactly the
  composites of Section~3, built from \(W_\phi\) instead of \(W\).
\item The random parts of the fit enter only through \(R\) and \(\Phi\),
  which are products of two small quantities.
\end{itemize}

\begin{figure}[h]\centering
\includegraphics[width=\textwidth]{T6_error_vs_leading.pdf}
\caption{\textbf{Checking the leading term point by point.}
Each dot is one time point in one of 20 simulated data sets (Gaussian
noise).  Horizontal axis: the leading term \(L^{(k)}_t\).  Vertical axis:
the actual error of component \(k\).  A dot on the diagonal means the
leading term predicts the error exactly.  The four panels are stages 1, 2,
5 and 8.
\emph{Takeaway:} the dots sit on the diagonal; the printed medians show the
typical miss is 30 to 150 times smaller than the typical error.}
\end{figure}

\section{Why the leading term is small at every stage}

The leading term is a sum over observations, \(L^{(k)}_t=\sum_u\xi_{t,u}(\varepsilon_u)\),
with independent pieces.  Each piece is bounded because the score
\(\rho'\) is bounded (\(|\rho'|\le1\)), and its size is set by the transfer
operators and the weights.  Bernstein's inequality then gives, for every
stage and time point at once,
\[
  |L^{(k)}_t|\lesssim\frac{C_W}{\lambda}\sqrt{\frac{z_x}{nh_k}} .
\]
\begin{itemize}
\item \(C_W\) is the sup-norm bound on the composites (Section~3, with bent
  weights), and it does not depend on \(k\).
\item Earlier stages have wider windows, so their weights are geometrically
  smaller; their contributions sum to a constant multiple of the last one.
\item Only the boundedness of \(\rho'\) is used, so this step holds under
  contamination as well.  That is the robust advantage.
\end{itemize}

The bent-weight composites are still bounded.  This is proved when the total
bending is small (\(C_P\tau_\Sigma\le\tfrac12\); the test signal gives 0.49).
In simulation it holds far beyond that: \(\le1\) even for the \(3\times\)
signal.

\section{The remainders, and where the picture breaks}

\begin{figure}[h]\centering
\includegraphics[width=\textwidth]{{{artifact:art_4a665027-9fe4-4ae1-a30b-74df4fcb47ee}}}
\caption{\textbf{Error budget by stage.}
Each line is the largest value over time, taken as a median over simulated
data sets, on a log scale.  Black: the error of component \(k\).  Orange: how
far that error is from the leading term.  Purple and green: the two parts of
the orange difference.  Purple is the remainder of the Fisher expansion at
that stage; green is the error of treating the effect of earlier errors as
linear.
The four panels are: Gaussian noise; the signal three times larger; and 20\%
of the points contaminated with spikes of size \(2\sigma\) or \(6\sigma\).
\emph{Takeaway:} orange stays well below black, except at the finest
stages under strong contamination, and there the cause is purple, the
remainder of a single stage.}
\end{figure}

\begin{itemize}
\item Without contamination the gap is 1--10\% of the error, and the Fisher
  remainder sets it.  Its size follows \(1/(nh_k)\), the second-order rate of
  Proposition~4.1.
\item With 20\% contamination at \(6\sigma\) the gap reaches 42\% at the
  31-point window.  There the single-pass remainder is \(0.128\) against an
  error of \(0.225\).  This is the local-sample-size condition of
  Proposition~4.1 failing at a small window; the propagation error is only
  \(0.038\).
\item Each remainder is of order \(z_x/(nh_j)\), and these sizes form a
  geometric series, so their sum is dominated by the last stage.
\end{itemize}

\section{What is proved and what is open}

\begin{center}\small
\begin{tabular}{p{0.55\textwidth}p{0.35\textwidth}}
\toprule
statement & status\\
\midrule
Exact decomposition \(e=L+\text{remainder}\) & proved\\
Linear composites bounded in sup norm, uniformly in \(k\), including Epanechnikov & proved (Fourier argument)\\
Bent-weight composites bounded when the bending is small & proved; holds numerically far beyond\\
Leading term \(\lesssim\sqrt{z_x/(nh_k)}\) uniformly, also under contamination & proved (Bernstein)\\
Fisher remainder at each stage & Proposition~4.1, applied stage by stage\\
Propagation remainder \(\Phi\) is second order & sketch; one U-statistic step open\\
Uniform rate \(A_1\sqrt{z_x/(nh_k)}+A_2z_x/(nh_k)\) & conditional on the line above\\
IMF.pdf true IMFs \(S^{(k)}\) vs \(\phi\)-path & equal numerically; proof open\\
Sharp constant, non-periodic design, finest windows under contamination & open\\
\bottomrule
\end{tabular}
\end{center}

\section{Cheat sheet}

\begin{center}\small
\begin{tabular}{ll}
\toprule
symbol & meaning (source)\\
\midrule
\(Y^{(k)},\ \IMFt^{(k)},\ h_k\) & residual, fitted IMF, bandwidth (IMF.pdf)\\
\(W^{(k)},\ A^{(k)}\) & linear smoother; product \((I-W^{(k)})\cdots(I-W^{(1)})\) (IMF.pdf)\\
\(S^{(k)}\) & true IMF, averaged over the actual residual noise (IMF.pdf \S3.2)\\
\(X^{(k)}_\phi,\ S^{(k)}_\phi\) & clean residual and target of the \(\phi\)-path (new)\\
\(F^{(k)}_t,\ \nabla^{(k)}_t,\ R^{(k)}\) & curvature, centered score, Fisher remainder at stage \(k\) (Prop.~4.1)\\
\(W^{(k)}_\phi\) & population equivalent kernel: linear weights bent by \(\phi''\) (new)\\
\(T^{(k,j)}\) & how a stage-\(j\) error reaches component \(k\) (new)\\
\(z_x\) & \(x+\log(4nK)\): union over all time points and stages\\
\bottomrule
\end{tabular}
\end{center}

\begin{idea}
\textbf{In five sentences.}  Each IRMF stage is a first-component problem
plus the errors handed down from earlier stages.  The Fisher expansion turns
each stage into a bent-weight linear smoother acting on bounded scores, plus
a second-order remainder.  The handed-down errors pass through chains of
\(I-W_\phi\), which in the sup norm stay bounded because the stages work at
well-separated scales.  Bounded scores give Bernstein bounds at every stage,
with the rate of the stage's own window.  What is left is second order and
small, except at the finest windows under heavy contamination.
\end{idea}

\end{document}
